% ---------------------------------------------------------------------
\chapter{Pull Requests und Code Review}
\label{kap:pr}

\begin{lernziele}
\begin{itemize}[nosep]
  \item verstehen, was ein Pull Request ist und warum er in der Industrie zentral ist,
  \item wissen, was bei einem Code Review geprüft wird,
  \item die drei Arten der Übernahme kennen und unterscheiden,
  \item Pull Requests mit \cmd{gh} erstellen und übernehmen.
\end{itemize}
\end{lernziele}

\section{Was ein Pull Request ist}

\begin{definition}{Pull Request}
Ein \begriff{Pull Request} (kurz PR) ist ein Antrag auf GitHub: „Bitte übernimm die Commits
meines Branches in \cmd{main}.“ Er bündelt die Änderungen, ihre Beschreibung, die
Diskussion darüber und die Ergebnisse automatischer Tests an einem Ort.
\end{definition}

Git selbst kennt keine Pull Requests. Sie sind eine Funktion der Plattform. Bei GitLab
heisst dasselbe \emph{Merge Request}.

Die Seite eines Pull Requests auf GitHub hat vier Reiter:

\begin{description}[style=nextline, leftmargin=1.2cm]
  \item[Conversation] Beschreibung, Kommentare und Freigaben.
  \item[Commits] die einzelnen Commits des Branches.
  \item[Checks] die Ergebnisse automatischer Tests.
  \item[Files changed] alle Änderungen als Diff. Hier findet das Review statt: Zu jeder
    Zeile kann man einen Kommentar schreiben.
\end{description}

\section{Code Review}

Beim \begriff{Code Review} liest mindestens eine zweite Person die Änderungen, bevor sie
übernommen werden. Typische Fragen dabei:

\begin{itemize}
  \item \textbf{Korrektheit:} Tut der Code, was er soll? Sind Sonderfälle bedacht, etwa
        leere Listen oder fehlende Eingaben?
  \item \textbf{Verständlichkeit:} Versteht jemand anderes den Code ohne Erklärung?
        Sind Namen sprechend?
  \item \textbf{Umfang:} Enthält der Pull Request nur, was zur Aufgabe gehört?
  \item \textbf{Sicherheit:} Stehen Passwörter im Code? Werden Eingaben geprüft?
  \item \textbf{Tests:} Ist die neue Funktion durch Tests abgesichert?
\end{itemize}

\begin{praxis}
Studien und Erfahrungen grosser Firmen zeigen, dass Reviews viele Fehler finden, bevor sie
Kundinnen und Kunden erreichen. Mindestens so wichtig ist aber ein Nebeneffekt: Wissen
verteilt sich im Team. Wenn nur eine Person einen Programmteil kennt und krank wird, steht
die Arbeit still. Durch Reviews kennen immer mindestens zwei Personen jede Änderung.

Darum gilt in Teams: \textbf{Kleine Pull Requests werden gründlich geprüft, grosse nur
überflogen.} Ein Pull Request sollte in 15 bis 30 Minuten verständlich sein.
\end{praxis}

\section{Automatische Prüfungen (CI)}

Bei jedem Pull Request kann GitHub automatisch Programme ausführen, zum Beispiel Tests,
Prüfungen des Codestils oder einen Testbuild. Das heisst \begriff{Continuous Integration}
(CI), auf GitHub umgesetzt mit \emph{GitHub Actions}. Die Anweisungen dafür liegen als
Datei im Repository unter \cmd{.github/workflows/}.

Das Ergebnis erscheint im Pull Request als grüner Haken oder rotes Kreuz. Rot bedeutet:
Etwas ist kaputt, die Übernahme wird blockiert. So gelangt fehlerhafter Code gar nicht erst
nach \cmd{main}.

\section{Drei Arten der Übernahme}

Beim Übernehmen eines Pull Requests bietet GitHub drei Varianten an. Ausgangslage: Der
Branch enthält drei Commits.

\begin{center}
\renewcommand{\arraystretch}{1.4}
\begin{tabularx}{\linewidth}{@{}>{\ttfamily}l>{\raggedright\arraybackslash}X>{\raggedright\arraybackslash}X@{}}
  \toprule
  \normalfont\textbf{Option} & \textbf{Ergebnis in main} & \textbf{Typischer Einsatz} \\
  \midrule
  --merge  & alle drei Commits bleiben erhalten, dazu ein \emph{Merge-Commit}, der die Übernahme festhält
           & wenn die einzelnen Schritte nachvollziehbar bleiben sollen; gut für den Einstieg \\
  --squash & die drei Commits werden zu \textbf{einem} zusammengefasst
           & sehr verbreitet in Firmen: \cmd{main} enthält genau einen Commit pro Aufgabe \\
  --rebase & die drei Commits werden einzeln an \cmd{main} angehängt, ohne Merge-Commit
           & wenn eine geradlinige Geschichte ohne Verzweigungen gewünscht ist \\
  \bottomrule
\end{tabularx}
\end{center}

\begin{erfahrung}
Es gibt keine „richtige“ Variante, nur eine einheitliche. Teams legen sich auf eine fest,
damit die Geschichte von \cmd{main} einheitlich aussieht. In diesem Handbuch wird
\cmd{--merge} verwendet, weil man im Log genau sieht, was wann übernommen wurde.
\end{erfahrung}

\section{Pull Requests auf der Kommandozeile}

\begin{lstlisting}[style=befehle]
gh pr create --fill
\end{lstlisting}

\begin{teile}
  pr create & Pull Request vom aktuellen Branch nach \cmd{main} erstellen \\
  --fill    & Titel und Beschreibung aus den Commit-Nachrichten übernehmen, ohne Rückfragen \\
\end{teile}

\begin{lstlisting}[style=befehle]
gh pr merge --merge --delete-branch
\end{lstlisting}

\begin{teile}
  pr merge        & den Pull Request des aktuellen Branches übernehmen \\
  --merge         & Art der Übernahme (siehe oben) \\
  --delete-branch & danach den Branch auf GitHub und lokal löschen und zu \cmd{main} wechseln \\
\end{teile}

Weitere nützliche Befehle:

\begin{lstlisting}[style=befehle]
gh pr list                             # offene Pull Requests anzeigen
gh pr view --web                       # aktuellen Pull Request im Browser öffnen
gh pr checks                           # Stand der automatischen Prüfungen
\end{lstlisting}

\section{Eine gute PR-Beschreibung}

Die Beschreibung richtet sich an die prüfende Person. Sie beantwortet drei Fragen:

\begin{lstlisting}[style=ausgabe]
## Was ändert sich?
Stack mit verketteter Liste statt Array implementiert.

## Warum?
Bei unbekannter Datenmenge braucht der Array-Stack teures Umkopieren.
Die Liste wächst in O(1) pro Element.

## Wie getestet?
Eigene Tests in stack_test.py, alle grün.
\end{lstlisting}

\begin{zusammenfassung}
\begin{itemize}[nosep]
  \item Ein Pull Request ist ein Antrag auf Übernahme mit Diskussion und Testergebnissen.
  \item Code Review findet Fehler und verteilt Wissen; kleine PRs werden besser geprüft.
  \item CI testet automatisch, rot blockiert die Übernahme.
  \item \cmd{--merge}, \cmd{--squash}, \cmd{--rebase}: Wichtig ist, dass ein Team sich auf eine Variante einigt.
\end{itemize}
\end{zusammenfassung}
